\documentclass[11pt]{amsart}
\usepackage[margin=1.05in]{geometry}
\usepackage{amsmath,amssymb,amsthm,mathtools,booktabs,hyperref,array}
\hypersetup{colorlinks=true,linkcolor=blue,citecolor=blue,urlcolor=blue}
\numberwithin{equation}{section}
\newtheorem{theorem}{Theorem}[section]
\newtheorem{lemma}[theorem]{Lemma}
\newtheorem{proposition}[theorem]{Proposition}
\theoremstyle{definition}
\newtheorem{definition}[theorem]{Definition}
\theoremstyle{remark}
\newtheorem{remark}[theorem]{Remark}
\newcommand{\Ftwo}{\mathbb F_2}
\newcommand{\tr}{\operatorname{Tr}}
\newcommand{\dimf}{\dim_{\Ftwo}}
\title[The Carlet zero-sum subspace conjecture]{A Proof and Lean Formalization of the Carlet Zero-Sum Subspace Conjecture}
\author{Wang Quanzhi}
\thanks{Email: \href{mailto:mrhwrn120504@gmail.com}{mrhwrn120504@gmail.com}.}
\author{ChatGPT}
\thanks{ChatGPT is an artificial intelligence system developed by OpenAI and contributed as a research-writing assistant. The human author accepts responsibility for the mathematical statements, code, and final manuscript.}
\date{November 2026}

\begin{document}
\begin{abstract}
Let $K$ be a finite field of cardinality $2^n$, viewed as an $n$-dimensional vector space over $\Ftwo$. We prove the full odd- and even-degree ranges of the Carlet zero-sum subspace conjecture: if $n$ is odd with $3\le k$ and $k+3\le n$, or $n$ is even with $2\le k$ and $k+2\le n$, then there is a $k$-dimensional $\Ftwo$-linear subspace $E\subseteq K$ such that
\[
  \sum_{x\in E\setminus\{0\}}x^{-1}=0.
\]
Both parity ranges have been formalized in Lean 4 with universal quantification over the corresponding finite fields. The odd-degree proof uses root-product polynomials, Hasse bounds, low-dimensional constructions, and energy iteration; the even-degree proof uses two- and three-dimensional base cases, dimension recurrence, and complementary-dimension symmetry. The complete default build, independent verification, and final axiom audit pass: there are no compiler warnings, source and OLean hashes match, and the audited dependencies are only \texttt{propext}, \texttt{Classical.choice}, and \texttt{Quot.sound}, with no \texttt{sorryAx} or custom axiom.
\end{abstract}

\keywords{finite fields, multiplicative inverse, zero-sum subspace, subspace polynomial, Kloosterman sum, Hasse bound, Lean formalization}
\subjclass[2020]{11T06, 12E20, 03B35, 68V20}

\maketitle

\section{Introduction and main theorem}

\noindent\textbf{Relationship to earlier research.} Cheng's October 2026 preprint~\cite{Cheng2026} already proves Carlet's existence conjecture and obtains a quantitative counting theorem. Our contribution is the explicitly documented Lean 4 formalization and its machine-checked constructive components and axiom audits; we do not claim the first proof of the existence conjecture.


Let $K$ be a finite field of characteristic $2$. Its nonzero elements form a finite group, but the inverse map $x\mapsto x^{-1}$ is not additive. This makes inverse sums on additive subspaces a natural finite-field problem: for which dimensions must a subspace contain a balanced collection of inverses? We study the zero-sum condition arising in Carlet's work~\cite{Carlet2025}, uniformly over all finite fields of a fixed cardinality.

For an $\Ftwo$-linear subspace $E\subseteq K$, write
\[
 S(E)=\sum_{x\in E\setminus\{0\}}x^{-1}.
\]
We state the odd-degree range first, give the even-degree range separately, and then combine them into a parity-dependent conclusion. The Lean formalization preserves the original quantifiers in both cases.

\begin{definition}
A finite-dimensional $\Ftwo$-linear subspace $E\subseteq K$ is a \emph{zero-sum subspace} if
\[
 S(E):=\sum_{x\in E\setminus\{0\}}x^{-1}=0.
\]
\end{definition}

\begin{theorem}[Odd-degree Carlet statement]\label{thm:main}
Let $n,k\in\mathbb N$, with $n$ odd, $3\le k$, and $k+3\le n$. For every finite field $K$ of characteristic $2$ and cardinality $2^n$, there exists a $k$-dimensional $\Ftwo$-linear subspace $E\subseteq K$ such that $S(E)=0$.
\end{theorem}

The assertion is uniform in $K$: the formal theorem quantifies over every field structure, finiteness structure, and $\Ftwo$-algebra structure satisfying the cardinality equality. The lower-range condition $k+3\le n$ is retained exactly. The even-degree assertion is stated separately below with its own base dimension and admissibility condition.

\section{Root products and a coefficient criterion}

For a finite $\Ftwo$-subspace $E\subseteq K$, set
\[
 P_E(T)=\prod_{u\in E}(T+u).
\]
This is a linearized polynomial. If $r=\dimf E$, then
\[
 P_E(T)=c_0T+c_1T^2+c_2T^4+\cdots+T^{2^r},\qquad c_0\ne0.
\]

\begin{lemma}[Reciprocal-sum criterion]\label{lem:coeff}
For every nonzero finite-dimensional $\Ftwo$-subspace $E$,
\[
 S(E)=\frac{c_1}{c_0},\qquad\text{and hence}\qquad
 S(E)=0\Longleftrightarrow [T^2]P_E=0.
\]
\end{lemma}
\begin{proof}
Separate the zero root:
\[
 P_E(T)=T\prod_{u\in E\setminus\{0\}}(T+u).
\]
The constant coefficient of the product on the right is $\prod_{u\ne0}u=c_0$. Its linear coefficient is
\[
 \sum_{u\ne0}\prod_{v\ne0,u}v
 =\left(\prod_{v\ne0}v\right)\left(\sum_{u\ne0}u^{-1}\right)=c_0S(E).
\]
After multiplication by $T$, this becomes the $T^2$ coefficient of $P_E$. Since $c_0\ne0$, the equivalence follows.
\end{proof}

Let $I(E)=P_E(K)$ and let $D(E)=I(E)^\perp$ with respect to the trace pairing. Then $\dimf I(E)=n-\dimf E$ and $\dimf D(E)=\dimf E$. If $\dimf E=k$ and $a_{k-1}=[T^{2^{k-1}}]P_E$, the formalized dual-polynomial identity gives
\begin{equation}\label{eq:duality}
 S(D(E))=a_{k-1}^{2}.
\end{equation}
Thus a $k$-dimensional space whose root product has vanishing next-to-leading Frobenius coefficient yields a $k$-dimensional zero-sum space. In general this zero-sum space is $D(E)$, not necessarily $E$ itself.

If $x\notin E$, grouping the roots into the two cosets of $E$ and using $P_E(T+x)=P_E(T)+P_E(x)$ gives the extension identity
\begin{equation}\label{eq:extension}
 P_{E+\langle x\rangle}(T)=P_E(T)^2+P_E(x)P_E(T).
\end{equation}
This formula drives the low-dimensional constructions.

\begin{lemma}[Complementary dimensions]\label{lem:complement}
If $\dimf K=n$ and $|K|>2$, then $K$ has a $k$-dimensional zero-sum subspace if and only if it has an $(n-k)$-dimensional zero-sum subspace.
\end{lemma}
\begin{proof}
The linear map $P_E:K\to K$ has kernel $E$ and image $I(E)$ of dimension $n-k$. The subspace polynomials satisfy
\[
 P_{I(E)}(P_E(T))=T^q-T.
\]
Write $c_0=[T]P_E\ne0$, $c_1=[T^2]P_E$, $d_0=[T]P_{I(E)}\ne0$, and $d_1=[T^2]P_{I(E)}$. Comparing $T^2$ coefficients gives $d_0c_1+d_1c_0^2=0$, so $c_1=0$ if and only if $d_1=0$. Lemma~\ref{lem:coeff} now shows that $E$ is zero-sum exactly when $I(E)$ is. Since $\dimf I(E)=n-k$, this proves one direction; applying the same dimension relation once more gives the reverse implication.
\end{proof}

\section{The three-dimensional case: an elliptic curve}

Write $q=|K|=2^n$ and define the binary Kloosterman sum
\[
 \operatorname{Kl}_K=\sum_{v\in K^\times}(-1)^{\tr(v+v^{-1})}.
\]
Consider the characteristic-two Weierstrass curve
\[
 C:\quad y^2+xy=x^3+1.
\]
Its discriminant is $1$, so it is elliptic. Counting affine points by fibers over $x$ and using the fact that $z^2+z=c$ has $1+(-1)^{\tr(c)}$ roots gives
\[
 \#C(K)=q+1+\operatorname{Kl}_K.
\]
The Hasse bound~\cite{Silverman} therefore yields $|\operatorname{Kl}_K|\le2\sqrt q$.

Partition the three-dimensional subspaces containing $1$ by their next Frobenius coefficient. The exact finite counting identity formalized in the project is
\begin{equation}\label{eq:p3}
 (q-2)(q-4)\,p_3(0)=q+1-3\operatorname{Kl}_K,
\end{equation}
where $p_3(0)$ is the proportion of the corresponding finite family for which that coefficient vanishes. For odd $n\ge7$, we have $q\ge128$ and
\[
 q+1-3\operatorname{Kl}_K\ge q+1-6\sqrt q>0.
\]
Both factors on the left of \eqref{eq:p3} are positive, so $p_3(0)>0$. A three-dimensional subspace with vanishing next coefficient exists; the coefficient-duality theorem and Lemma~\ref{lem:coeff} turn it into a zero-sum subspace.

\begin{proposition}\label{prop:k3}
For every odd $n\ge7$, every finite field $K$ of cardinality $2^n$ contains a three-dimensional zero-sum subspace.
\end{proposition}

\section{The four-dimensional case: Fourier counting on a norm-one conic}

We prove the four-dimensional case for every admissible odd degree. The count must impose two trace-orthogonality conditions at once, so that two values lie in the image of a chosen subspace polynomial.

\subsection{A two-dimensional seed and the extension identity}

Choose $t\in K\setminus\{0,1\}$ and put $H=\langle1,t\rangle_{\Ftwo}$. Since $n$ is odd, $T^2+T+1$ has no root in $K$, so $c=t^2+t+1\ne0$. Direct multiplication gives
\[
 P_H(T)=T^4+cT^2+(t^2+t)T.
\]
Consider the norm-one conic
\begin{equation}\label{eq:normcurve}
 \mathcal C(K)=\{(a,u)\in K^2:a^2+au+u^2=1\}.
\end{equation}
For odd $n$, this conic has $q+1$ points. Define the linear map $L(x)=c^{-2}P_H(x)$ and let $D=(\operatorname{im}L)^\perp$, where orthogonality is with respect to the nondegenerate trace pairing $(r,s)\mapsto\tr(rs)$. Since $\ker L=H$, we have $\dimf\operatorname{im}L=n-2$, $\dimf D=2$, and $D^\perp=\operatorname{im}L$. Consequently, whenever $a,u\in D^\perp$, there are $x,y\in K$ such that
\[
 P_H(x)=c^2a,\qquad P_H(y)=c^2u.
\]

Impose on a point of $\mathcal C(K)$ the two conditions $u\in D^\perp$ and $a+u\in D^\perp$, equivalently
\[
 \tr(ru)=0,\quad \tr(s(a+u))=0\qquad(r,s\in D).
\]
Write $\chi(z)=(-1)^{\tr(z)}$. For a nonzero frequency pair $(A,B)$, the Fourier coefficient is
\[
 \mathcal A(A,B)=\sum_{x\in K}\chi\!\left(\frac{Bx^2+A}{x^2+x+1}\right)+\chi(B).
\]
It is obtained by counting the fibers of the Artin--Schreier equation
\[
 \delta^2+\delta=\frac{(A+B)x^2+Ax}{x^2+x+1}.
\]
A birational change of variables sends the generic solutions to the Weierstrass curve
\[
 y^2+xy+By=x^3+(A^2+B+1)x^2+(A+(A^2+B))x+B(A+(A^2+B))+A^2.
\]
Its discriminant is $(A^2+AB+B^2)^2$. Odd extension degree ensures that $T^2+T+1$ has no root in $K$ and that this discriminant is nonzero for every nonzero frequency pair. The Hasse bound controls the total curve-point count; the map from generic points to affine equation solutions is injective, and the exceptional fibers together account for at most five points. It follows that each nonzero-frequency coefficient satisfies $\mathcal A(A,B)\ge-2\sqrt q-4$. Using the four elements of $D$ as additive-character filters for each condition yields a 16-term Fourier expansion. The zero-frequency term is $q+1$, while the other 15 terms satisfy the bound above. Thus the number $N$ of points satisfying both conditions satisfies
\begin{equation}\label{eq:k4count}
 16N\ge q+1-30\sqrt q-60.
\end{equation}
If $q\ge2048$, the right-hand side is greater than $48$, so $N>3$.

For every point on the conic, $a=0$ forces $u=1$, while $a=1$ forces $u=0$ or $u=1$. There are at most three points with $a\in\{0,1\}$. Choose a good point with $a\ne0,1$. Since $u,a+u\in D^\perp$, also $a\in D^\perp$, so the preimages $x,y$ above exist.

Set $V=H+\langle x\rangle$. Then $\dimf V=3$, and \eqref{eq:extension} gives
\[
 [T^4]P_V=c^2+P_H(x)=c^2(1+a)\ne0.
\]
Furthermore,
\begin{align*}
 P_V(y)&=P_H(y)^2+P_H(x)P_H(y)\\
 &=c^4(u^2+au)=c^4(1+a^2)=\bigl(c^2(1+a)\bigr)^2.
\end{align*}
Here we used $a^2+au+u^2=1$ and $(1+a)^2=1+a^2$ in characteristic two. The displayed nonzero value implies $y\notin V$. Extending once more by $y$, identity \eqref{eq:extension} shows that the $T^8$ coefficient of the resulting four-dimensional root product vanishes. The same-dimension coefficient-duality theorem then gives a four-dimensional zero-sum subspace.

\begin{proposition}\label{prop:k4large}
For every odd $n\ge11$, every finite field $K$ of cardinality $2^n$ contains a four-dimensional zero-sum subspace.
\end{proposition}

The remaining admissible degrees $n=7,9,11,13,15$ are handled by exact finite-field constructions formalized in the project. Combining them with Proposition~\ref{prop:k4large} gives the full base family:
\begin{theorem}\label{thm:k4}
For every odd $n\ge7$ and every finite field $K$ of cardinality $2^n$, there is a four-dimensional zero-sum subspace.
\end{theorem}

\section{The five-dimensional case: a three-dimensional seed and a norm-one count}

\subsection{A nonzero-coefficient seed from dimension four}

By Theorem~\ref{thm:k4} and coefficient duality, choose a four-dimensional subspace $W$ with $[T^8]P_W=0$. Let $U$ be a three-dimensional hyperplane of $W$, and take $v\in W\setminus U$. Put $b=[T^4]P_U$. Comparing $T^8$ coefficients in \eqref{eq:extension} yields
\[
 0=[T^8]P_W=b^2+P_U(v).
\]
Since the roots of $P_U$ are exactly $U$, we have $P_U(v)\ne0$, and hence $b\ne0$. Put $c=b^2\ne0$.

Let $L=c^{-1}P_U$. Its image has codimension three. Set $D=(\operatorname{im}L)^\perp$; then $\dimf D=3$ and $|D|=8$. Impose on a point $(a,u)$ of \eqref{eq:normcurve} the conditions $u,a+u\in D^\perp$. Filtering both conditions over the eight elements of $D$ produces a 64-term Fourier expansion. The zero-frequency term is $q+1$, and each of the other 63 terms is at least $-2\sqrt q-4$. Therefore
\begin{equation}\label{eq:k5count}
 64N\ge q+1-63(2\sqrt q+4)=q+1-126\sqrt q-252.
\end{equation}
For $q\ge65536$, the right-hand side exceeds $192$, hence $N>3$. At most one point on the conic has $a=0$, and at most two have $a=1$. We may therefore choose a good point with $a\ne0,1$.

Because $u,a+u\in D^\perp=\operatorname{im}L$, their sum $a$ also belongs to the image. Choose $x,y\in K$ with
\[
 P_U(x)=ca,\qquad P_U(y)=cu.
\]
As $ca\ne0$, we have $x\notin U$. Let $V=U+\langle x\rangle$, so $\dimf V=4$. Write $A=[T^8]P_V$. Since $P_U$ is monic of degree eight, \eqref{eq:extension} gives
\[
 A=b^2+P_U(x)=c(1+a)\ne0.
\]
The conic equation implies $u^2+au=1+a^2$. Thus
\[
 P_V(y)=P_U(y)^2+P_U(x)P_U(y)
       =c^2(u^2+au)=c^2(1+a^2)=A^2\ne0.
\]
So $y\notin V$. With $Z=V+\langle y\rangle$, formula \eqref{eq:extension} yields
\[
 [T^{16}]P_Z=A^2+P_V(y)=0.
\]
The five-dimensional coefficient-duality theorem converts this vanishing into a five-dimensional zero-sum subspace. This conversion is essential: the $T^{16}$ coefficient is not itself the reciprocal sum; the proved same-dimension duality connects them.

\begin{proposition}\label{prop:k5large}
For every odd $n\ge17$, every finite field $K$ of cardinality $2^n$ contains a five-dimensional zero-sum subspace.
\end{proposition}

For the other admissible odd degrees $n=9,11,13,15$, exact finite-field constructions are formalized using explicit field models, linear-algebra witnesses, and a proof that the result transfers to every field of the same cardinality. Hence:
\begin{theorem}\label{thm:k5}
For every odd $n\ge9$ and every finite field $K$ of cardinality $2^n$, there is a five-dimensional zero-sum subspace.
\end{theorem}

\section{Large dimensions: an actual probability recurrence and energy decay}

Let $\Omega_d$ be the finite set of $d$-dimensional $\Ftwo$-subspaces containing $1$. For $E\in\Omega_d$, let $b(E)$ denote its next Frobenius coefficient, and define the exact distribution
\[
 p_d(b)=\frac{\#\{E\in\Omega_d:b(E)=b\}}{\#\Omega_d},\qquad
 f_d(b)=p_d(b)-q^{-1}.
\]
Then $\sum_{b\in K}f_d(b)=0$. The Lean development derives the recurrence for these distributions from the underlying finite-subspace counts; no abstract probabilistic model is substituted for the actual family.

For a nontrivial additive character $\psi:K\to\mathbb R$ and a real parameter $s$, define the transfer operator by
\[
 (T_sf)(b)=\frac{1}{q-s}\begin{cases}
 (q-s)f(0),&b=1,\\[2pt]
 \displaystyle\sum_{v\ne0}\left(1+(s-1)\psi\!\left(\frac{b-1}{v}\right)\right)f(v),&b\ne1.
 \end{cases}
\]
Equivalently, the special branch is at shifted coordinate $b-1=0$; this is the convention used in the implementation. Put $\mathcal E(f)=\sum_b f(b)^2$.

\begin{lemma}[Two-step energy contraction]\label{lem:energy}
If $s\ge2$ and $32s^2\le q$, then every zero-sum function $f$ satisfies
\[
 \mathcal E(T_{2s}(T_sf))\le\frac{4s^2}{q}\mathcal E(f).
\]
If also $s^2\le q$ and $s<q$, then one step is nonexpansive: $\mathcal E(T_sf)\le\mathcal E(f)$.
\end{lemma}
\begin{proof}
Decompose the transfer as a rank-one perturbation in the additive-character direction. Orthogonality of finite additive characters gives an exact squared-norm identity, including the exceptional coordinate term. For two successive ranks, the scalar on the distinguished direction is bounded by $2s/\sqrt q$; squaring gives $4s^2/q$. The condition $32s^2\le q$ makes the denominators $q-s$ and $q-2s$ positive and controls the remaining rank-one term. The same exact identity gives one-step nonexpansion when $s^2\le q$. The formal proof establishes the identity before applying these inequalities.
\end{proof}

The initial one-dimensional distribution is a point mass and has centered energy $1-q^{-1}$. Iterating Lemma~\ref{lem:energy} along odd dimensions gives, for $k=2r+1$,
\[
 \mathcal E(f_k)\le \frac{2^{2r^2+2r}}{q^r}(1-q^{-1}).
\]
For $k=2r+2$, apply one final nonexpansive step to obtain the same bound. If $p_k(0)=0$, then $f_k(0)=-q^{-1}$, and this coordinate alone implies $\mathcal E(f_k)\ge q^{-2}$. Thus a strict upper bound below $q^{-2}$ forces $p_k(0)>0$.

For odd $k=2r+1$, the hypotheses $r\ge5$ and $n\ge4r+3$ imply
\[
 2r^2+2r<n(r-2),
\]
so $2^{2r^2+2r}/q^r<q^{-2}$. For even $k=2r+2$, the same strict inequality follows from $r\ge4$ and $n\ge4r+5$. Therefore the coefficient-zero probability is positive for odd $k\ge11$ and even $k\ge10$, whenever $n\ge2k+1$. A positive probability means an actual subspace has the desired vanishing coefficient; coefficient duality then gives a zero-sum subspace.

\section{Even degree: base cases, recurrence, and complementary dimensions}
\sloppy

The even-degree argument starts from different base cases. If $n$ is even, $K$ contains its unique subfield $\mathbb F_4$. Choose $\omega\in\mathbb F_4\setminus\mathbb F_2$, so $\omega^2+\omega+1=0$. The sum of the inverses of its three nonzero elements is $1+\omega+\omega^2=0$, giving the $k=2$ base. For even $n\ge6$, a three-dimensional base follows from the Kloosterman sum and the Hasse bound for an elliptic curve.

\begin{theorem}[Even-degree statement]\label{thm:even}
Let $n$ be even and $|K|=2^n$. If $2\le k$ and $k+2\le n$, then $K$ contains a $k$-dimensional zero-sum subspace.
\end{theorem}
\begin{proof}
We first establish the three-dimensional base. Put $q=|K|$ and define $\operatorname{Kl}_K=\sum_{v\in K^\times}(-1)^{\operatorname{Tr}(v+v^{-1})}$. The point count for $y^2+xy=x^3+1$ gives $|\operatorname{Kl}_K|\le2\sqrt q$. Partitioning the three-dimensional subspaces containing $1$ by their next Frobenius coefficient yields the exact even-degree identity
\[
 (q-2)(q-4)p_3(0)=q-13+3\operatorname{Kl}_K.
\]
In even degree $\operatorname{Tr}(1)=0$, so the character-weighted count differs from the odd-degree formula. For $n\ge6$, we have $q\ge64$, and the Hasse bound gives
\[
q-13+3\operatorname{Kl}_K\ge q-13-6\sqrt q=(\sqrt q-3)^2-22\ge3>0.
\]
Since $(q-2)(q-4)>0$, it follows that $p_3(0)>0$. Thus a three-dimensional subspace with vanishing coefficient exists, and the coefficient criterion together with same-dimension duality gives a three-dimensional zero-sum subspace.

Now suppose $2k\le n$. Even dimensions start from the two-dimensional base, and each next case follows by adding a two-dimensional zero-sum subspace to one of dimension $k-2$. Odd dimensions start from the three-dimensional base just proved and use the same step. The dimension-addition lemma requires $2(k-2)<n$, which follows from $2k\le n$. This proves the result for every $2\le k\le n/2$.

If $2k>n$, put $k'=n-k$. The assumption $k+2\le n$ gives $k'\ge2$, and $2k'<n$. Apply the lower-half result to $k'$ and then use complementary-dimension symmetry to obtain dimension $k$. When $n=4$, the hypotheses allow only $k=2$, which is the base case.
\end{proof}

\section{Odd-degree dimension assembly}

The low-dimensional cases also generate dimensions $6,7,8,9$. Suppose $E,F$ are zero-sum subspaces of dimensions $k,\ell$ and $k\ell<n$. A formalized linear-algebra lemma gives a nonzero scalar $a$ with $aF\subseteq P_E(K)$. Define
\[
 G=P_E^{-1}(aF).
\]
Since $\ker P_E=E$ and $\dimf P_E(K)=n-k$, we get $\dimf G=k+\ell$. The composition formula for subspace polynomials and the coefficient criterion show that $G$ is zero-sum. Thus
\begin{equation}\label{eq:add}
 \text{zero-sum dimensions }k,\ell\text{ and }k\ell<n
 \quad\Longrightarrow\quad k+\ell\text{ is a zero-sum dimension}.
\end{equation}
Using dimensions $3$ and $4$, the pairs $(3,3),(3,4),(4,4),(3,6)$ give dimensions $6,7,8,9$. Dimensions $k\ge10$ are covered by the energy argument. Hence for the lower half $2k+1\le n$, the only base dimensions required are $3,4,5$.

For the upper half $2k+1>n$, put $k'=n-k$. Since $n$ is odd and $k+3\le n$, one checks that $k'\ge3$ and $2k'+1\le n$. Apply the lower-half result to $k'$, then use Lemma~\ref{lem:complement} to obtain the desired $k$-dimensional space. This accounts for every admissible odd-degree parameter.

\begin{proof}[Proof of Theorem~\ref{thm:main}]
If $2k+1\le n$, use Proposition~\ref{prop:k3}, Theorem~\ref{thm:k4}, or Theorem~\ref{thm:k5} for $k=3,4,5$ respectively; use \eqref{eq:add} for $6\le k\le9$; and use the energy argument for $k\ge10$. If $2k+1>n$, reduce to $k'=n-k$ and apply complementary-dimension symmetry. The finite-field constructions and all transitions work for every $K$ of cardinality $2^n$, so the original universal quantifier is preserved.
\end{proof}

\begin{theorem}[Combined odd- and even-degree statement]\label{thm:full}
Let $n,k\in\mathbb N$, and let $K$ be a finite field of characteristic $2$ and cardinality $2^n$. If
\[
 (\operatorname{Odd}(n)\land 3\le k\land k+3\le n)
 \quad\text{or}\quad
 (\operatorname{Even}(n)\land 2\le k\land k+2\le n),
\]
then $K$ contains a $k$-dimensional zero-sum subspace.
\end{theorem}
\begin{proof}
The odd case is Theorem~\ref{thm:main}; the even case is Theorem~\ref{thm:even}. Both hold uniformly for every finite field of the specified cardinality.
\end{proof}

\section{Formal verification and reproducibility}\label{sec:formal}
\sloppy

The formalization uses Lean 4 and mathlib~\cite{Lean4,Mathlib}. The original odd-degree statement is defined as the proposition
\[
\forall n,k,\quad \operatorname{Odd}(n)\to 3\le k\to k+3\le n\to
\forall K,\quad |K|=2^n\to \operatorname{HasZeroSumSubspace}(K,k).
\]
{\raggedright
The original statement is defined in \texttt{Statement.lean}; the odd-degree theorem is \texttt{oddCarletConjecture\_main} in \texttt{MainTheorem.lean}.
The unconditional even theorem, \texttt{evenDegreeCarlet\_unconditional}, is in \texttt{EvenDegree.lean}. It covers every even $n$, every finite field $K$ with $|K|=2^n$, and every $k$ with $2\le k$ and $k+2\le n$.
Together these declarations imply Theorem~\ref{thm:full}. The main proof entry points are:
\par}
\begin{center}
\small\begin{tabular}{>{\raggedright\arraybackslash}p{0.43\textwidth}>{\raggedright\arraybackslash}p{0.47\textwidth}}
\toprule
Proof component & Lean source\\
\midrule
Reciprocal sum and root-product criterion & \texttt{Elementary.lean}, \texttt{Statement.lean}\\
Three-dimensional Hasse/Kloosterman argument & \texttt{KloostermanThree.lean}\\
All odd-degree four-dimensional cases & \texttt{K4AllOddDegreesProbe.lean} and imports\\
Five-dimensional norm-one construction & \texttt{K5NormOneConstructionProbe.lean} and imports\\
Distribution recurrence and energy contraction & \texttt{ProbabilityEnergy.lean}, \texttt{TwoStepEnergy.lean}\\
Dimension addition and complementary reduction & \texttt{AdditiveClosure.lean}, \texttt{LowDimensionReduction.lean}\\
Original quantified theorem & \texttt{MainTheorem.lean}\\
Even-degree bases and dimension recurrence & \texttt{EvenDegree.lean}, \texttt{KloostermanThree.lean}\\
Combined odd/even parameter ranges & \texttt{MainTheorem.lean}, \texttt{EvenDegree.lean}\\
\bottomrule
\end{tabular}
\end{center}

{\raggedright
The 11 November 2026 acceptance record uses Lean 4.33.0. The complete default build checked 274 modules; the independent verifier audited 942 declarations and checked 188 Hasse dependency modules. There were zero compiler warnings, and all source/OLean hashes match.
Both main theorems passed independent statement checks. Their axiom sets contain only \texttt{propext}, \texttt{Classical.choice}, and \texttt{Quot.sound}; the audit found no \texttt{sorryAx} or custom unproved axiom.
The Hasse bound is imported from the fixed AINTLIB snapshot~\cite{AINTLIB} (\texttt{1c1c7466\allowbreak 4e40071c\allowbreak 2c2165bc\allowbreak 55ca2616\allowbreak a67ccd6b}); its source is \texttt{HasseBound.lean} and it is licensed under Apache 2.0.
The pinned mathlib source baseline is \texttt{db584cd6\allowbreak d46c92f2\allowbreak 09a44c0f\allowbreak 1c829460\allowbreak d327499d}. Reproduction commands are:
\par}
\begin{verbatim}
cd formalization
python -I build.py
python -I verify_build.py
python -I audit_final_theorem.py
\end{verbatim}
The full-build summary and final axiom audit are recorded in \texttt{verified\_build\_summary.json} and \texttt{final\_theorem\_audit.json}.

The odd-degree proof follows the route checked by Lean: subspace polynomials, Hasse-controlled character sums, norm-one counts, and energy iteration. The even-degree proof uses two- and three-dimensional bases, dimension addition, and complementary-dimension symmetry. Finite computations are used only in the explicitly bounded small-degree branches, with the relevant constructions and transfer to arbitrary fields of the same cardinality checked in Lean.

\section*{AI-use statement}
ChatGPT assisted with research organization, formalization support, and manuscript drafting. The human author is responsible for the mathematical content, verification record, and final text.

\begin{thebibliography}{9}
\bibitem{Carlet2025}
C.~Carlet, \emph{On the vector subspaces of $\mathbb F_{2^n}$ over which the multiplicative inverse function sums to zero}, Designs, Codes and Cryptography 93 (2025), 1237--1254. DOI: \href{https://doi.org/10.1007/s10623-024-01531-6}{10.1007/s10623-024-01531-6}.

\bibitem{Cheng2026}
K.~Cheng, \emph{Counting zero-sum subspaces for the multiplicative inverse function}, arXiv:2610.06445 (2026), \url{https://arxiv.org/abs/2610.06445}.
\bibitem{Silverman}
J.~H.~Silverman, \emph{The Arithmetic of Elliptic Curves}, 2nd ed., Graduate Texts in Mathematics 106, Springer, 2009.

\bibitem{Lean4}
L.~de Moura and S.~Ullrich, \emph{The Lean 4 Theorem Prover and Programming Language}, in: Automated Deduction -- CADE 28, Lecture Notes in Computer Science 12699, Springer, 2021, pp.~625--635.

\bibitem{Mathlib}
The mathlib Community, \emph{The Lean Mathematical Library}, Proceedings of the 9th ACM SIGPLAN International Conference on Certified Programs and Proofs (CPP 2020), pp.~367--381. DOI: \href{https://doi.org/10.1145/3372885.3373824}{10.1145/3372885.3373824}.

\bibitem{AINTLIB}
C.~Birkbeck, \emph{AINTLIB}, Hasse--Weil bound module, project snapshot \texttt{1c1c7466\allowbreak 4e40071c\allowbreak 2c2165bc\allowbreak 55ca2616\allowbreak a67ccd6b} (Apache License 2.0).
\end{thebibliography}
\end{document}
